\section{Dinámica de resistencia entre capas}

El análisis empírico revela una dinámica de \textbf{filtrado probabilístico bidireccional} entre niveles de integración de la persistencia. «Resistencia» se entiende como un \textbf{efecto estructural de filtrado probabilístico} emergente de la configuración de acoplamiento entre componentes. No se postula ninguna fuerza activa ni agente. Se describe una variación observable y medible en la probabilidad de que una intensificación local genere reorganización estructural según el nivel de antisincronización del sistema.

\subsection{«Resistencia ascendente»}

Cuando la antisincronización es elevada, la probabilidad de que una intensificación local (Capa 1) logre modificar la matriz de coherencias del sistema (Capa 3) disminuye sustancialmente. Los módulos mantienen escalas de persistencia $\tau_s$ dispersas; la fragmentación resultante reduce la permeabilidad relacional necesaria para que los cambios locales se propaguen. En términos estrictamente operacionales, el aumento de la desviación estándar de $\tau_s$ entre componentes actúa como un filtro que eleva el umbral requerido para que una intensificación local se traduzca en reorganización estructural detectable mediante la norma de Frobenius.

\subsection{«Resistencia descendente»}

Cuando la antisincronización ha descendido (Capa 2 consolidada), la probabilidad de que intensificaciones locales aún fragmentadas alteren la estructura global ya alcanzada también se reduce. Una organización coherente de escalas de persistencia genera una configuración en la que las fluctuaciones locales que no alcanzan un umbral mínimo de reducción de fragmentación tienden a ser absorbidas sin modificar la matriz de relaciones global. Operacionalmente, la baja antisincronización estabiliza la estructura global frente a intensificaciones locales que no logran alinear suficientemente las escalas de persistencia de los módulos.

\subsection{Implicación ontológica}

La antisincronización funciona como regulador de la permeabilidad entre capas: su valor modula la probabilidad de transición de la Capa 1 a la Capa 3. El presente de nivel superior (Capa 3) no se reduce a la suma de presentes locales intensificados; emerge como una estructura relacional cuya estabilidad depende de la reducción previa de fragmentación. Las transiciones entre capas siguen siendo los eventos discretos ontológicamente relevantes del tiempo extramental detectable mediante RECD. Las denominaciones «resistencia ascendente» y «resistencia descendente» se mantienen entre comillas para indicar que se trata de efectos de filtrado probabilístico en direcciones opuestas de la jerarquía, y no de procesos activos de protección o defensa por parte de ningún agente o subsistema.